% Section 5.3, from lectures/L05/README.md: what you should be able to do afterwards, the
% questions to test yourself, and the next lecture.

\section{Review}
\label{c5:sec:review}

\needspace{6\baselineskip}
\subsection*{What you should be able to do now}
After this chapter, you should be able to:
\begin{itemize}
  \item Say what a major and a minor number each select, and find both for any node in \file{/dev}.
  \item Write a \code{struct file_operations} and register it, both the long way and with
    \code{misc_register}.
  \item Explain the three separate things \code{copy_to_user} checks, and what goes wrong if you
    skip it.
  \item Say what \code{read()} must return when it has nothing to give, and why zero is the wrong
    answer.
  \item Choose the right \code{-Exxx} for a given failure and say what userspace will print for it.
  \item Give one reason to add an \code{ioctl} and two reasons not to.
  \item Write a ring buffer that survives its own wraparound, and prove it with a test rather than
    by inspection.
\end{itemize}

\needspace{6\baselineskip}
\subsection*{Questions to test yourself}
\begin{enumerate}
  \item Two device nodes have the same major and different minors. What do they have in common?
  \item Why is \code{copy_to_user} a function rather than a macro that expands to \code{memcpy}?
  \item Your \code{read()} returns 0 whenever the buffer is empty. Describe what \code{cat
    /dev/yours} does, and say why it is not what you wanted.
  \item What does \code{open()} on your device return if you never implemented \code{open} in your
    \code{file_operations}?
  \item Where does the node in \file{/dev} come from, given that your driver only registered a
    number?
  \item A ring buffer passes every test except one about wraparound. What is almost certainly wrong
    with it, and why did the other tests not catch it?
\end{enumerate}

\needspace{6\baselineskip}
\subsection*{Where to look}
\begin{itemize}
  \item \file{lectures/L05/lab/qa_fifo.h} is the contract, printed in \secref{c5:app:b1}, and
    \file{lectures/L05/lab/qa_fifo_kunit.c} is what checks it. Reading the suite before you start is
    allowed and recommended; it is a specification written in a second language.
  \item The suite tests nothing about concurrency, and says so at the top. That is
    Chapter~\ref{c7:ch}.
\end{itemize}

\needspace{6\baselineskip}
\subsection*{Looking ahead}
Chapter~\ref{c6:ch} is about:
\begin{itemize}
  \item Where a driver gets memory, and the flag that says which context you are in.
  \item How a driver reaches its hardware, and why the compiler must not be trusted with it.
  \item \code{ioremap}, \code{readl}, and the address you compute by hand before you map it.
  \item The error path you write four times and then delete.
\end{itemize}
